\chapter{Mesoscopic Model of Excitatory Clustering}
\label{ch:mesoscopic-excitatory-clustering}

\textit{A clustered spiking network becomes mathematically useful when each cluster can be treated as a noisy population with its own activity variable. This chapter builds the minimal mesoscopic description for excitatory clustering: recurrent self-excitation creates up-states, shared inhibition creates competition, and finite cluster size turns deterministic attractors into metastable switching.}

% ============================================================
\section{Treating Each Cluster as a Population}
\label{sec:treating-each-cluster-as-population}
% ============================================================

Start with $Q$ excitatory clusters. Cluster $\alpha$ contains a set of excitatory neurons
%
\begin{equation}
  \mathcal{C}_{\alpha} \subset \{1,\ldots,N_E\},
  \qquad
  n_{\alpha} = |\mathcal{C}_{\alpha}|,
  \qquad
  \sum_{\alpha=1}^{Q} n_{\alpha} = N_E .
  \label{eq:cluster-sets}
\end{equation}
%
The index $\alpha$ labels a mesoscopic unit, not a single neuron. The size $n_{\alpha}$ controls how much intrinsic population noise remains after averaging over neurons. Large clusters have smoother activity; small clusters switch more easily.

For neuron $i$, write its spike train as
%
\begin{equation}
  x_i(t)=\sum_k \delta(t-t_i^k).
  \label{eq:neuron-spike-train-ch10}
\end{equation}
%
The raw population activity of cluster $\alpha$ is the spike train average
%
\begin{equation}
  A_{\alpha}(t)
  =
  \frac{1}{n_{\alpha}}
  \sum_{i\in \mathcal{C}_{\alpha}} x_i(t).
  \label{eq:raw-cluster-activity}
\end{equation}
%
$A_{\alpha}(t)$ has units of spikes per neuron per unit time. It measures the instantaneous output of the cluster, but it is too singular for ordinary differential equations because it is a sum of impulses.

The filtered cluster rate $r_{\alpha}(t)$ is a smoother state variable:
%
\begin{keyequation}[Cluster Rate as Filtered Population Activity]
\begin{equation}
  \tau_s \frac{\dd r_{\alpha}}{\dd t}
  =
  -r_{\alpha} + A_{\alpha}(t).
  \label{eq:filtered-cluster-rate}
\end{equation}
\tcblower
$A_{\alpha}$ is the raw spike output of the cluster, $r_{\alpha}$ is the synaptically filtered activity seen by downstream neurons, and $\tau_s$ is the filtering time scale.
\end{keyequation}

\noindent Increasing $r_{\alpha}$ means that cluster $\alpha$ is contributing more recurrent excitation and more drive to inhibition. Increasing $n_{\alpha}$ at fixed mean rate does not change the deterministic mean, but it reduces the relative fluctuations around that mean.

This is the central mesoscopic move: keep the cluster identity, discard individual neuron identities inside the cluster, and represent the cluster by a small number of population variables. The learner should check whether observables of interest, such as cluster lifetimes, number of active clusters, and slow autocorrelation, can be computed from $r_{\alpha}(t)$ without reconstructing every spike.

% ============================================================
\section{Cluster-Level Rate Equations}
\label{sec:cluster-level-rate-equations}
% ============================================================

A deterministic mesoscopic rate model has the form
%
\begin{equation}
  \tau_E \dot r_{\alpha}
  =
  -r_{\alpha}
  +
  \Phi_E(u_{\alpha}),
  \qquad
  \alpha=1,\ldots,Q .
  \label{eq:e-cluster-rate-equation}
\end{equation}
%
The transfer function $\Phi_E$ maps total input $u_{\alpha}$ to the firing rate that the cluster would express if the input were held fixed. It is a mesoscopic surrogate for the single-neuron spiking model plus synaptic bombardment. Its slope is a gain: high slope means small input changes produce large rate changes.

The input to cluster $\alpha$ is
%
\begin{equation}
  u_{\alpha}
  =
  h_{\alpha}
  +
  \sum_{\beta=1}^{Q} J_{\alpha\beta} r_{\beta}
  -
  J_{EI} r_I .
  \label{eq:e-cluster-input}
\end{equation}
%
Here $h_{\alpha}$ is external drive to the cluster, $J_{\alpha\beta}$ is effective excitation from cluster $\beta$ to cluster $\alpha$, $r_I$ is the activity of the inhibitory pool, and $J_{EI}>0$ is the strength of inhibitory input onto excitatory clusters.

The simplest inhibitory pool is one shared population:
%
\begin{equation}
  \tau_I \dot r_I
  =
  -r_I
  +
  \Phi_I\!\left(
    h_I + J_{IE}\sum_{\beta=1}^{Q} q_{\beta} r_{\beta}
  \right),
  \qquad
  q_{\beta}=\frac{n_{\beta}}{N_E}.
  \label{eq:shared-inhibition}
\end{equation}
%
The weights $q_{\beta}$ convert cluster rates into the mean excitatory activity seen by the inhibitory pool. If all clusters have equal size, then $q_{\beta}=1/Q$.

Equations \eqref{eq:e-cluster-rate-equation}--\eqref{eq:shared-inhibition} are not yet a microscopic derivation. They are a controlled bookkeeping layer: each term has a direct circuit meaning, and each parameter can be matched to a spiking simulation or varied in a phase diagram.

% ============================================================
\section{Effective Recurrent Excitation Within Clusters}
\label{sec:effective-recurrent-excitation-within-clusters}
% ============================================================

Excitatory clustering means that neurons inside the same cluster excite each other more strongly, or more frequently, than neurons in different clusters. At the mesoscopic level this becomes a structured coupling matrix:
%
\begin{equation}
  J_{\alpha\beta}
  =
  \begin{cases}
    J_{EE}^{+}, & \alpha=\beta,\\
    J_{EE}^{-}, & \alpha\ne\beta.
  \end{cases}
  \label{eq:e-cluster-coupling}
\end{equation}
%
$J_{EE}^{+}$ is the within-cluster recurrent excitation. $J_{EE}^{-}$ is the between-cluster excitation. The difference $J_{EE}^{+}-J_{EE}^{-}$ is the strength of clustering.

Often the total expected excitatory input is held approximately fixed while clustering is varied. For equal-sized clusters, a convenient normalization is
%
\begin{equation}
  p_{\mathrm{in}} J_{EE}^{+}
  +
  (1-p_{\mathrm{in}})J_{EE}^{-}
  =
  J_{EE}^{0},
  \label{eq:fixed-mean-excitation}
\end{equation}
%
where $p_{\mathrm{in}}$ is the fraction of an excitatory neuron's recurrent excitatory input that comes from its own cluster, and $J_{EE}^{0}$ is the baseline unclustered coupling. Solving for the between-cluster coupling gives
%
\begin{equation}
  J_{EE}^{-}
  =
  \frac{J_{EE}^{0}-p_{\mathrm{in}}J_{EE}^{+}}{1-p_{\mathrm{in}}}.
  \label{eq:jminus-normalization}
\end{equation}
%
Increasing $J_{EE}^{+}$ while enforcing \eqref{eq:jminus-normalization} concentrates excitation inside clusters and weakens excitation across clusters.

The effect on dynamics is immediate. A cluster that becomes slightly more active receives extra self-excitation through $J_{EE}^{+}r_{\alpha}$. If the transfer gain $\Phi_E'(u_{\alpha})$ is large enough, that small increase can be amplified into an up-state. The learner should compute how the fixed points of a single cluster change as $J_{EE}^{+}$ rises while the mean input constraint is kept fixed.

% ============================================================
\section{Inter-Cluster Inhibition and Competition}
\label{sec:inter-cluster-inhibition-competition}
% ============================================================

In an excitatory-clustered network, clusters do not usually inhibit each other directly. They compete indirectly through the shared inhibitory pool. When cluster $\beta$ goes up, it drives $r_I$ upward; the increased inhibition then suppresses every excitatory cluster, including cluster $\alpha$.

This can be made explicit by eliminating fast inhibition. Suppose inhibition is faster than the excitatory cluster dynamics and linearize the inhibitory transfer around an operating point:
%
\begin{equation}
  r_I
  \approx
  r_I^0
  +
  G_I J_{IE}
  \sum_{\beta=1}^{Q} q_{\beta}(r_{\beta}-r_{\beta}^0),
  \qquad
  G_I=\Phi_I'(u_I^0).
  \label{eq:linearized-fast-inhibition}
\end{equation}
%
$G_I$ is the inhibitory gain. It tells us how strongly the inhibitory population reacts to a change in excitatory drive.

Substituting \eqref{eq:linearized-fast-inhibition} into \eqref{eq:e-cluster-input} gives an effective excitatory coupling
%
\begin{keyequation}[Effective Competition Matrix]
\begin{equation}
  K_{\alpha\beta}
  =
  J_{\alpha\beta}
  -
  J_{EI}G_I J_{IE}q_{\beta}.
  \label{eq:effective-competition-matrix}
\end{equation}
\tcblower
$J_{\alpha\beta}$ is direct recurrent excitation, while $J_{EI}G_I J_{IE}q_{\beta}$ is the indirect inhibitory feedback caused by activity in cluster $\beta$.
\end{keyequation}

\noindent Off-diagonal competition is strong when the indirect inhibitory term exceeds the between-cluster excitation:
%
\begin{equation}
  J_{EI}G_I J_{IE}q_{\beta} > J_{EE}^{-}.
  \label{eq:competition-condition}
\end{equation}
%
In that regime, raising $r_{\beta}$ decreases the net input to cluster $\alpha\ne\beta$. This is the mesoscopic origin of winner-take-more dynamics: clusters excite themselves but inhibit competitors through a shared inhibitory pathway.

% ============================================================
\section{Fixed Points: No Active, One Active, and Many Active Clusters}
\label{sec:fixed-points-cluster-count}
% ============================================================

A fixed point is a vector $(r_1^*,\ldots,r_Q^*,r_I^*)$ satisfying $\dot r_{\alpha}=0$ and $\dot r_I=0$. Fixed points are candidate network states. In clustered networks, the most useful classification is by the number $k$ of active clusters.

Assume equal cluster sizes and symmetric coupling. A $k$-active state has
%
\begin{equation}
  r_{\alpha}^*
  =
  \begin{cases}
    r_U, & \alpha \in \mathcal{A},\\
    r_D, & \alpha \notin \mathcal{A},
  \end{cases}
  \qquad
  |\mathcal{A}|=k ,
  \label{eq:k-active-ansatz}
\end{equation}
%
where $r_U$ is the up-state rate and $r_D$ is the down-state rate. The set $\mathcal{A}$ identifies which clusters are active; symmetry means that only its size $k$ matters for the fixed-point equations.

The fixed-point equations are
%
\begin{align}
  r_U
  &=
  \Phi_E\!\left(
    h
    +
    J_{EE}^{+}r_U
    +
    (k-1)J_{EE}^{-}r_U
    +
    (Q-k)J_{EE}^{-}r_D
    -
    J_{EI}r_I
  \right),
  \label{eq:k-active-up}\\
  r_D
  &=
  \Phi_E\!\left(
    h
    +
    kJ_{EE}^{-}r_U
    +
    J_{EE}^{+}r_D
    +
    (Q-k-1)J_{EE}^{-}r_D
    -
    J_{EI}r_I
  \right),
  \label{eq:k-active-down}\\
  r_I
  &=
  \Phi_I\!\left(
    h_I + \frac{J_{IE}}{Q}\bigl(kr_U+(Q-k)r_D\bigr)
  \right).
  \label{eq:k-active-inhibition}
\end{align}
%
The first equation says an active cluster must be self-consistent as an up-state. The second says inactive clusters must remain down despite excitation from active clusters. The third says inhibition tracks the average excitatory activity.

Three cases organize the phase diagram:
%
\begin{enumerate}[leftmargin=*]
  \item \textbf{No active clusters} ($k=0$): all clusters sit near the low-rate state $r_D$. This can represent spontaneous asynchronous activity before a cluster escapes upward.
  \item \textbf{One active cluster} ($k=1$): one cluster is up and suppresses the others. This is the cleanest attractor-selection state.
  \item \textbf{Many active clusters} ($k>1$): several clusters are co-active. This can represent a multi-item state, a broad metastable state, or an unstable intermediate depending on parameters.
\end{enumerate}

The practical computation is to solve \eqref{eq:k-active-up}--\eqref{eq:k-active-inhibition} for each $k$ and then test stability. The number of mathematically possible fixed points can be large because every subset $\mathcal{A}$ of size $k$ is a symmetry-related state.

% ============================================================
\section{Stability of Cluster Up-States}
\label{sec:stability-of-cluster-upstates}
% ============================================================

Fixed points only matter dynamically if perturbations behave predictably near them. Linearize the excitatory subsystem around a fixed point. With fast inhibition already folded into $K_{\alpha\beta}$, the Jacobian is
%
\begin{equation}
  \mathcal{J}_{\alpha\beta}
  =
  \frac{1}{\tau_E}
  \left[
    -\delta_{\alpha\beta}
    +
    \Phi_E'(u_{\alpha}^*)K_{\alpha\beta}
  \right].
  \label{eq:e-cluster-jacobian}
\end{equation}
%
$-\delta_{\alpha\beta}/\tau_E$ is the leak of the rate variable. The term $\Phi_E'(u_{\alpha}^*)K_{\alpha\beta}/\tau_E$ is recurrent feedback: coupling multiplied by gain. A fixed point is linearly stable if every eigenvalue of $\mathcal{J}$ has negative real part.

For a single isolated cluster coordinate, the scalar stability condition is
%
\begin{equation}
  \Phi_E'(u^*)K_{\mathrm{self}} < 1 .
  \label{eq:scalar-upstate-stability}
\end{equation}
%
This inequality can look surprising: strong self-excitation is needed to create the up-state, but at the stable up-state the local feedback slope must be less than the leak. Bistability is created by nonlinear shape, not by making the stable point locally explosive.

For a symmetric $k$-active state, two perturbation types are especially informative. Let $U$ denote the active block and $D$ denote the inactive block. In the contrast-mode shortcut below,
%
\begin{equation}
  K_{UU}=K_{\alpha\alpha}\ \text{for}\ \alpha\in\mathcal{A},
  \qquad
  K_{UB}=K_{\alpha\beta}\ \text{for}\ \alpha,\beta\in\mathcal{A},\ \alpha\ne\beta,
  \label{eq:active-block-couplings}
\end{equation}
%
and
%
\begin{equation}
  K_{DD}=K_{\alpha\alpha}\ \text{for}\ \alpha\notin\mathcal{A},
  \qquad
  K_{DB}=K_{\alpha\beta}\ \text{for}\ \alpha,\beta\notin\mathcal{A},\ \alpha\ne\beta.
  \label{eq:inactive-block-couplings}
\end{equation}
%
Here the second subscript $B$ means ``between two clusters inside the same activity block,'' not a new population. Thus $K_{UU}$ and $K_{DD}$ are diagonal effective self-couplings, while $K_{UB}$ and $K_{DB}$ are off-diagonal effective couplings within the active and inactive blocks. For equal-size clusters with the competition matrix in \eqref{eq:effective-competition-matrix}, these reduce to the corresponding diagonal and off-diagonal entries of $K_{\alpha\beta}$.
%
\begin{align}
  \lambda_{\mathrm{active},c}
  &\approx
  \frac{-1+\Phi_E'(u_U)(K_{UU}-K_{UB})}{\tau_E},
  \label{eq:active-contrast-eigenvalue}\\
  \lambda_{\mathrm{inactive},c}
  &\approx
  \frac{-1+\Phi_E'(u_D)(K_{DD}-K_{DB})}{\tau_E}.
  \label{eq:inactive-contrast-eigenvalue}
\end{align}
%
The active-contrast mode raises one active cluster while lowering another. If it is unstable, equal co-activation breaks into a more selective state. The inactive-contrast mode tests whether a down cluster can separate from the inactive pool and begin an escape.

The learner should not only ask whether an up-state exists. The sharper question is which perturbation destroys it: amplitude decay, competition among active clusters, or escape of an inactive cluster.

% ============================================================
\section{Bistability at the Mesoscopic Level}
\label{sec:bistability-at-mesoscopic-level}
% ============================================================

Bistability means that the same input supports both a low-rate and a high-rate stable state. A single cluster embedded in a slowly varying background can be reduced to
%
\begin{equation}
  \tau_E \dot r
  =
  -r
  +
  \Phi_E\!\left(h_{\mathrm{eff}}+J_{\mathrm{eff}}r\right),
  \label{eq:single-cluster-bistability}
\end{equation}
%
where $h_{\mathrm{eff}}$ includes external drive, background excitation, and baseline inhibition, while $J_{\mathrm{eff}}$ is the net recurrent feedback of the cluster after inhibition is accounted for.

Define
%
\begin{equation}
  G(r)
  =
  \Phi_E(h_{\mathrm{eff}}+J_{\mathrm{eff}}r)-r .
  \label{eq:bistability-nullcline-function}
\end{equation}
%
Fixed points are zeros of $G(r)$. A stable fixed point has $G'(r)<0$; an unstable fixed point has $G'(r)>0$. A bistable cluster has three zeros:
%
\begin{equation}
  r_D < r_S < r_U ,
  \label{eq:down-saddle-up}
\end{equation}
%
where $r_D$ is the down-state, $r_U$ is the up-state, and $r_S$ is the saddle threshold separating their basins in the one-dimensional reduction.

At a saddle-node bifurcation, two fixed points merge. The defining conditions are
%
\begin{equation}
  G(r^*)=0,
  \qquad
  G'(r^*)=\Phi_E'(h_{\mathrm{eff}}+J_{\mathrm{eff}}r^*)J_{\mathrm{eff}}-1=0 .
  \label{eq:ch10-saddle-node-conditions}
\end{equation}
%
The first condition says the rate is self-consistent. The second says recurrent gain exactly balances leak. These equations locate the boundary of the bistable region in a phase diagram.

At the network level, cluster bistability becomes combinatorial multistability. If many subsets of clusters can be active, the deterministic system has many attractors or long-lived states. Finite-size fluctuations then move the network between them.

% ============================================================
\section{Finite-Size Noise and Metastable Transitions}
\label{sec:finite-size-noise-metastable-transitions}
% ============================================================

The deterministic equation gives the mean drift of a large cluster. A finite cluster has noisy spike counts. In a short window $\Delta t$, a near-Poisson approximation gives
%
\begin{equation}
  \E[\Delta N_{\alpha}]
  \approx
  n_{\alpha}r_{\alpha}\Delta t,
  \qquad
  \Var[\Delta N_{\alpha}]
  \approx
  n_{\alpha}r_{\alpha}\Delta t .
  \label{eq:cluster-count-moments}
\end{equation}
%
After dividing by $n_{\alpha}$, the relative size of the fluctuation scales like $1/\sqrt{n_{\alpha}}$. This is the basic finite-size law: larger clusters are less noisy.

A Langevin-level mesoscopic model writes
%
\begin{keyequation}[Finite-Size Cluster Dynamics]
\begin{equation}
  \dd r_{\alpha}
  =
  f_{\alpha}(\vec r)\,\dd t
  +
  \frac{1}{\sqrt{n_{\alpha}}}
  \sum_{\mu} B_{\alpha\mu}(\vec r)\,\dd W_{\mu}(t),
  \label{eq:cluster-langevin}
\end{equation}
\tcblower
$f_{\alpha}$ is deterministic drift, $B_{\alpha\mu}$ sets noise amplitude and correlations, $W_{\mu}$ are independent Wiener processes, and $n_{\alpha}^{-1/2}$ is the finite-size scaling.
\end{keyequation}

\noindent In the simplest diagonal approximation, $B_{\alpha\alpha}^2(\vec r)$ is proportional to the instantaneous activity of cluster $\alpha$. More refined mesoscopic theories include refractoriness, synaptic filtering, and shared input, which make $B$ state-dependent and non-diagonal.

Metastable switching is noise-driven escape from one basin of attraction to another. For a one-dimensional reduction with an effective potential $U(r)$, the mean escape time has the qualitative Kramers form
%
\begin{equation}
  T_{\mathrm{escape}}
  \approx
  C
  \exp\!\left(n_{\alpha}\Delta U\right),
  \label{eq:kramers-cluster-scaling}
\end{equation}
%
where $\Delta U$ is the barrier between the stable state and the saddle. The exact prefactor and barrier are model-dependent, but the scaling is the key point: if switching is finite-size driven, cluster lifetime should grow rapidly with cluster size.

% ============================================================
\section{Reproducing the E-Clustered Baseline}
\label{sec:reproducing-e-clustered-baseline}
% ============================================================

A mesoscopic E-clustered model should not merely show interesting dynamics. It should reproduce the specific observables that made the microscopic clustered network useful.

The minimal target list is:
%
\begin{enumerate}[leftmargin=*]
  \item \textbf{Cluster up- and down-states}: the distribution of $r_{\alpha}(t)$ should be bimodal or at least strongly skewed toward low and high activity regimes.
  \item \textbf{Slow cluster autocorrelation}: the autocorrelation of $r_{\alpha}(t)$ should decay on the scale of cluster lifetimes, not only synaptic time constants.
  \item \textbf{Occupancy distribution}: the model should match the distribution of the number of active clusters $k(t)$.
  \item \textbf{Switching statistics}: cluster lifetime distributions should be comparable to microscopic simulations at the same effective cluster size.
  \item \textbf{Stimulus bias}: external input to a subset of clusters should change occupancy and lifetimes in the expected direction.
\end{enumerate}

The occupancy variable is
%
\begin{equation}
  k(t)
  =
  \sum_{\alpha=1}^{Q}
  \mathbf{1}\{r_{\alpha}(t)>\theta_{\mathrm{up}}\},
  \label{eq:occupancy-variable}
\end{equation}
%
where $\theta_{\mathrm{up}}$ is a threshold chosen between down- and up-state rates. $k(t)$ is a coarse but powerful observable because it compresses the high-dimensional state into the number of active attractor-like clusters.

The mesoscopic model is successful only if its parameters have interpretable roles. $J_{EE}^{+}$ should deepen up-states, shared inhibition should control competition, $n_{\alpha}$ should control switching noise, and stimulus drive $h_{\alpha}$ should bias the occupancy distribution.

% ============================================================
\section{Minimal Simulation Tasks for Reproduction}
\label{sec:minimal-simulation-tasks-reproduction}
% ============================================================

The most efficient reproduction plan starts with the mesoscopic equations before returning to microscopic simulations.

First, sweep $J_{EE}^{+}$ and external drive $h$ in the deterministic system. For each point, solve the fixed-point equations and classify the stable states by $k$. This produces a deterministic phase diagram with asynchronous, bistable, and multi-active regions.

Second, add finite-size noise using \eqref{eq:cluster-langevin}. At fixed deterministic parameters, vary $n_{\alpha}$. If switching is finite-size driven, the mean lifetime should increase strongly with $n_{\alpha}$ while the locations of deterministic fixed points remain almost unchanged.

Third, compare mesoscopic and microscopic observables using the same definitions. For example,
%
\begin{equation}
  C_{\alpha}(\tau)
  =
  \frac{
    \E[(r_{\alpha}(t)-\bar r_{\alpha})(r_{\alpha}(t+\tau)-\bar r_{\alpha})]
  }{
    \Var[r_{\alpha}(t)]
  }
  \label{eq:ch10-cluster-autocorrelation}
\end{equation}
%
measures the normalized cluster autocorrelation. Its decay time should track cluster switching, not merely the synaptic filter.

Fourth, add a stimulus as a structured drive
%
\begin{equation}
  h_{\alpha}(t)
  =
  h_0 + s(t)m_{\alpha},
  \label{eq:structured-stimulus-drive}
\end{equation}
%
where $m_{\alpha}$ marks stimulus-targeted clusters. The computation is to compare $P(k)$, lifetimes, and variability before and after stimulus onset.

This minimal program separates three questions: where attractors exist, how finite-size noise moves between them, and how stimuli reshape the occupancy landscape.

\begin{chaptersummary}

\begin{center}
\renewcommand{\arraystretch}{1.45}
\begin{tabularx}{\textwidth}{@{}l X X@{}}
  \toprule
  \textbf{Concept} & \textbf{Equation} & \textbf{Key Insight} \\
  \midrule
  Cluster activity &
  $A_{\alpha}=n_{\alpha}^{-1}\sum_{i\in\mathcal{C}_{\alpha}}x_i$ &
  Average spikes inside one cluster \\
  Filtered rate &
  $\tau_s\dot r_{\alpha}=-r_{\alpha}+A_{\alpha}$ &
  Synaptic-scale population variable \\
  E-cluster dynamics &
  $\tau_E\dot r_{\alpha}=-r_{\alpha}+\Phi_E(u_{\alpha})$ &
  Leak plus nonlinear input-output map \\
  Clustered excitation &
  $J_{\alpha\alpha}=J_{EE}^{+}$, $J_{\alpha\beta}=J_{EE}^{-}$ &
  Self-excitation creates up-state candidates \\
  Competition &
  $K_{\alpha\beta}=J_{\alpha\beta}-J_{EI}G_I J_{IE}q_{\beta}$ &
  Shared inhibition turns activity in one cluster into suppression of others \\
  Bistability &
  $G(r)=0$, $G'(r)=0$ at saddle-node &
  Up- and down-states appear through nonlinear self-consistency \\
  Finite-size switching &
  $\dd r=f\,\dd t+n^{-1/2}B\,\dd W$ &
  Cluster size controls stochastic escape \\
  \bottomrule
\end{tabularx}
\end{center}

\end{chaptersummary}

\begin{conceptquestions}
  \item What information is retained when a cluster is replaced by $r_{\alpha}(t)$? What information about individual neurons is discarded?

  \item Why does increasing $J_{EE}^{+}$ while preserving mean recurrent input make clusters more likely to have up-states?

  \item In the effective coupling $K_{\alpha\beta}=J_{\alpha\beta}-J_{EI}G_I J_{IE}q_{\beta}$, which term is responsible for competition? How would increasing inhibitory gain change the phase diagram?

  \item For a $k$-active fixed point, why must both the active-rate equation and inactive-rate equation be checked?

  \item Explain why stable up-states require local feedback slope less than leak even though strong recurrent excitation is needed to create the up-state.

  \item If cluster lifetimes do not increase when $n_{\alpha}$ is increased, what does that suggest about the switching mechanism?
\end{conceptquestions}
